\documentclass[aspectratio=169]{beamer}

\usepackage{xcolor}

% Dark background
\setbeamercolor{background canvas}{bg=black}

\usetheme{Madrid}
\definecolor{primary}{RGB}{64,60,104}
\definecolor{secondary}{RGB}{94,234,212}

\definecolor{bgcolor}{RGB}{42,44,52}
\definecolor{textcolor}{RGB}{232,228,220}
\definecolor{accent}{RGB}{125,115,245}

\setbeamercolor{background canvas}{bg=bgcolor}
\setbeamercolor{normal text}{fg=textcolor,bg=bgcolor}
\setbeamercolor{structure}{fg=accent}

\title{ARM64 HPC Research Topics}
\subtitle{Graduate Thesis Topics in HPC and Data Analytics}
\author{Dr. John Burns}
\date{\today}

\begin{document}

\begin{frame}
\titlepage
\end{frame}

%====================================================

\begin{frame}{1. Energy Efficiency of ARM64 for Scientific HPC Workloads}
\small

\textbf{Research Question}\\
Under which workload characteristics does ARM64 deliver better performance-per-watt than x86-64 for scientific HPC applications?

\vspace{0.3em}

\textbf{Abstract}\\
This research evaluates whether modern ARM64 processors provide better energy efficiency than x86-64 systems for scientific HPC workloads. The study benchmarks numerical and parallel applications to analyze runtime, power consumption, thermal behavior, and scalability.

\vspace{0.3em}

\textbf{Tooling}\\
LINPACK, HPCG, Linux perf, ARM Streamline, OpenMP

\vspace{0.3em}

\textbf{References}
\begin{itemize}
\item F. Mantovani et al. ``Performance and Energy Consumption of HPC Workloads on a Cluster Based on ARM ThunderX2 CPUs.'' \emph{FGCS}, 2020.

\item N. Simakov et al. ``Are We Ready for Broader Adoption of ARM in the HPC Community?'' \emph{HPCAsia Workshops}, 2023.
\end{itemize}

\end{frame}

%====================================================

\begin{frame}{2. Memory-Centric Profiling on ARM HPC Systems}
\small

\textbf{Research Question}\\
Can ARM SPE-based memory profiling identify optimization opportunities that traditional CPU profiling misses?

\vspace{0.35em}

\textbf{Abstract}\\
This thesis investigates whether memory-centric profiling techniques can identify bottlenecks and improve optimization strategies for ARM-based HPC systems. The research focuses on cache utilization, memory locality, bandwidth saturation, and profiling-guided optimization for scientific applications.

\vspace{0.35em}

\textbf{Tooling}\\
perf, ARM SPE, PAPI, STREAM benchmark

\vspace{0.35em}

\textbf{References}
\begin{itemize}
\item Samuel Miksits, Johannes Weidendorfer, and Dieter Kranzlmuller. ``Multi-level Memory-Centric Profiling on ARM Processors with ARM Statistical Profiling Extension (SPE).'' \emph{arXiv preprint arXiv:2410.01514}, 2024.

\item Arm Ltd. ``Statistical Profiling Extension.'' Arm Developer Documentation and Architecture Blog Series, 2023--2024.
\end{itemize}

\end{frame}

%====================================================

\begin{frame}{3. Transformer Inference on ARM Many-Core CPUs}
\small

\textbf{Research Question}\\
What software and hardware factors determine transformer inference throughput and latency on ARM many-core CPUs?

\vspace{0.35em}

\textbf{Abstract}\\
This research investigates transformer inference performance on ARM many-core CPUs. The study analyzes thread scheduling, memory usage, vectorization efficiency, quantization, and runtime optimization techniques for modern AI workloads.

\vspace{0.35em}

\textbf{Tooling}\\
PyTorch, ONNX Runtime, Hugging Face Transformers, perf

\vspace{0.35em}

\textbf{References}
\begin{itemize}
\item Canqun Jiang, Zheng Wang, Xiaoyang Wang, and David Kaeli. ``Full-Stack Optimizing Transformer Inference on ARM Many-Core CPUs.'' \emph{IEEE Transactions on Parallel and Distributed Systems}, vol. 34, no. 9, pp. 2478--2491, 2023.

\item Canqun Jiang, Zheng Wang, and David Kaeli. ``Characterizing and Optimizing Transformer Inference on ARM Many-Core Processors.'' \emph{Proceedings of the International Workshop on Performance, Portability and Productivity in HPC}, 2022.
\end{itemize}

\end{frame}

%====================================================

\begin{frame}{4. AI-Assisted Autotuning for ARM HPC Workloads}
\small

\textbf{Research Question}\\
Can machine-learning-based autotuning improve performance and energy efficiency on ARM HPC systems?

\vspace{0.35em}

\textbf{Abstract}\\
This thesis explores the use of machine learning and Bayesian optimization techniques for automatic tuning of HPC workloads on ARM architectures. The research investigates parameter optimization, compiler tuning, scheduling decisions, and performance prediction for scientific applications.

\vspace{0.35em}

\textbf{Tooling}\\
OpenTuner, GPTune, Python, LLVM/Clang, perf

\vspace{0.35em}

\textbf{References}
\begin{itemize}
\item Yurii Oleynik, Jack Deslippe, and Samuel Williams. ``GPTune: Multitask Learning for Autotuning HPC Applications.'' \emph{IEEE Transactions on Parallel and Distributed Systems}, vol. 32, no. 9, pp. 2229--2245, 2021.

\item Nikolay Simakov, Hartwig Anzt, and Thomas Papatheodore. ``Are We Ready for Broader Adoption of ARM in the HPC Community?'' \emph{Proceedings of the ACM International Conference on High Performance Computing in Asia-Pacific Workshops}, 2023.
\end{itemize}

\end{frame}

%====================================================

\begin{frame}{5. Memory-Bound Parallel Workloads on ARM64 and x86-64}
\small

\textbf{Research Question}\\
How do ARM64 and x86-64 architectures differ in cache efficiency, memory bandwidth saturation, and OpenMP scaling for memory-bound scientific workloads?

\vspace{0.35em}

\textbf{Abstract}\\
This thesis investigates memory-bound scientific applications on ARM64 and x86-64 systems. The study evaluates cache efficiency, memory bandwidth utilization, NUMA behavior, and OpenMP scalability to identify architectural differences that influence performance and energy efficiency.

\vspace{0.35em}

\textbf{Tooling}\\
STREAM, OpenMP, Linux perf, PAPI, LIKWID

\vspace{0.35em}

\textbf{References}
\begin{itemize}
\item Filippo Mantovani, Edoardo Calore, Alberto Schenck, and Thomas Fahringer. ``Performance and Energy Consumption of HPC Workloads on a Cluster Based on ARM ThunderX2 CPUs.'' \emph{Future Generation Computer Systems}, vol. 103, pp. 126--138, 2020.

\item Zhen Shi, Hartwig Anzt, and Jack Dongarra. ``ARM SVE Unleashed: Performance and Insights Across HPC Applications on NVIDIA Grace.'' \emph{arXiv preprint arXiv:2505.09462}, 2025.
\end{itemize}

\end{frame}

%====================================================

\begin{frame}{Conclusion}
\small

\textbf{Key Research Themes}
\begin{itemize}
\item Energy-efficient high performance computing
\item Memory-centric optimization and profiling
\item AI inference on ARM many-core processors
\item Machine-learning-assisted performance tuning
\item Memory-bound scientific application analysis
\end{itemize}

\vspace{0.5em}

\textbf{Why ARM64?}
\begin{itemize}
\item Rapid growth in HPC and cloud infrastructure
\item Improved performance-per-watt characteristics
\item Increasing adoption in AI and scientific computing
\item Strong opportunities for systems analytics research
\end{itemize}

\end{frame}

\end{document}

